\section{Discussion}
\label{sec:discussion}

The decade of work surveyed above suggests that SE is best understood not as a
single clustering score, but as a unifying principle for graph intelligence:
systems become intelligible when their uncertainty can be reduced by a
compact hierarchy.
This interpretation explains why SE appears in apparently distant settings,
including graph pooling, chromatin-domain discovery, event detection, social
bot analysis, and reinforcement learning.
In each case, the method searches for a structural abstraction that preserves
the dominant organization while suppressing incidental variation.

For TGINA's focus on graph intelligence and network applications, SE is
especially useful because it connects three concerns that are often handled
separately.
First, it is a \emph{measure}: the entropy value summarizes how ordered or
disordered a network is under a candidate hierarchy.
Second, it is an \emph{optimizer}: minimizing the measure constructs the
hierarchy itself.
Third, it is a \emph{learning signal}: differentiable variants allow the same
principle to shape GNN encoders, pooling layers, and decision abstractions.
This measure--optimizer--learning-signal triad is the core reason SE deserves
a full journal treatment rather than a short method catalogue.

At the same time, the benchmark results caution against overclaiming.
Topology-only SE objectives are most appropriate when communities correspond
to structural bottlenecks.
On attributed citation networks, semantic labels may be better explained by
text features than by citation topology, and feature-augmented models such as
LSENet can therefore dominate in accuracy.
Conversely, on noisy synthetic topologies, differentiable SE remains
competitive with or stronger than standard baselines.
The practical lesson is that SE should be matched to the source of structure:
pure topology for latent abstraction and network organization, feature-aware
SE for semantic graph learning, and constrained SE for domain-specific
structures such as genomic intervals.

Finally, SE offers a promising language for trustworthy graph intelligence in
LLM and autonomous-agent systems.
GraphRAG pipelines, tool-use traces, memory graphs, and multi-agent interaction
networks all produce evolving relational structures whose reliability depends
on whether useful organization can be separated from noisy connectivity.
SE can quantify that organization, detect disorder, and potentially guide the
construction of stable retrieval or decision hierarchies.
This connection remains early, but it gives a concrete path from classical
structural information theory to the emerging graph intelligence problems
that motivate the TGINA venue.
